\section{Self-Supervised Representation Learning}
\label{sec:ssl}

Data augmentation (\cref{subsec:data-augmentation}) reduces
overfitting by diversifying the encoder's inputs. But augmentation
does not change what the encoder is trained to predict -- the
Q-learning loss remains the sole objective shaping the encoder's
representations. In the 100K regime, this signal is weak: Q-targets
are bootstrapped from the agent's own poor estimates, and the
encoder may learn shallow features that satisfy current targets
without capturing the full structure of the environment.
Self-supervised representation learning addresses this gap by adding
an auxiliary training objective that shapes the encoder's
representations independently of the reward signal. This section
introduces the general idea, the key technical challenge
(representational collapse), and the specific method used in our
experiments: Self-Predictive Representations
(SPR)~\citep{schwarzer2021dataefficient}.

\subsection{Shaping Better Representations}
\label{subsec:shaping-better-representations}

In the DQN architecture (\cref{subsec:dqn-architecture}), raw frames
pass through three convolutional layers and a fully connected layer,
producing a 512-dimensional representation of the game state.
Everything the agent ``knows'' about the current observation passes
through this vector. The Q-values for all actions are computed from
it, and the Q-learning loss (\cref{eq:dqn-loss}) is the only training
signal that shapes it. The encoder learns to extract whatever features
predict Q-values, and nothing more (\cref{fig:single-objective}).

\begin{figure}[htbp]
\centering
% Thesis-wide categorical color palette
% Usage: % Thesis-wide categorical color palette
% Usage: \input{figures/palette} in any figure file
%
% Muted primary colors -- print-friendly, visually distinct.
% Inspired by Cooper Union's red/blue primaries, softened for
% academic figures. Use palA-palC for data series in scatter
% plots, line charts, bar charts, etc.

\definecolor{palA}{HTML}{1F77B4}   % Tableau blue
\definecolor{palB}{HTML}{D62728}   % Tableau red
\definecolor{palC}{HTML}{2CA02C}   % Tableau green
 in any figure file
%
% Muted primary colors -- print-friendly, visually distinct.
% Inspired by Cooper Union's red/blue primaries, softened for
% academic figures. Use palA-palC for data series in scatter
% plots, line charts, bar charts, etc.

\definecolor{palA}{HTML}{1F77B4}   % Tableau blue
\definecolor{palB}{HTML}{D62728}   % Tableau red
\definecolor{palC}{HTML}{2CA02C}   % Tableau green

\definecolor{encfade}{HTML}{B8D2E6}
\definecolor{outfront}{HTML}{DCE8F4}
\definecolor{outtop}{HTML}{EBF1F8}
\definecolor{outside}{HTML}{C2D6E8}
\definecolor{fcfront}{HTML}{B3CEE4}
\definecolor{fctop}{HTML}{CDDEF0}
\definecolor{fcside}{HTML}{8DB3D2}

\newcommand{\cuboid}[7]{%
    \fill[#5, draw=black!35, line join=round]
        (#1, {-#3/2}, 0) -- ({#1+#2}, {-#3/2}, 0) --
        ({#1+#2}, {#3/2}, 0) -- (#1, {#3/2}, 0) -- cycle;
    \fill[#6, draw=black!35, line join=round]
        (#1, {#3/2}, 0) -- ({#1+#2}, {#3/2}, 0) --
        ({#1+#2}, {#3/2}, #4) -- (#1, {#3/2}, #4) -- cycle;
    \fill[#7, draw=black!35, line join=round]
        ({#1+#2}, {-#3/2}, 0) -- ({#1+#2}, {#3/2}, 0) --
        ({#1+#2}, {#3/2}, #4) -- ({#1+#2}, {-#3/2}, #4) -- cycle;
}

\begin{tikzpicture}[
    x={(1cm,0cm)}, y={(0cm,1cm)}, z={(0.35cm,0.25cm)},
    arr/.style={->, thick, >=stealth, black!45},
    gradarr/.style={<-, thick, >={Stealth[length=6pt, width=5pt]}, dashed, palA!55},
    lbl/.style={font=\scriptsize, text=black!45},
]

% Observation label
\node[font=\small, text=black!85] at (-0.5, 0, 0) {$s_t$};
\draw[arr] (0.0, 0, 0) -- (0.6, 0, 0);

% Encoder (trapezoid -- smaller)
\fill[encfade, draw=black!60, thick, line join=round]
    (0.8, -0.65, 0) -- (0.8, 0.65, 0) -- (2.2, 0.4, 0)
    -- (2.2, -0.4, 0) -- cycle;
\node[font=\small, text=black!80] at (1.4, 0, 0) {$f_\theta$};

% z_t label
\node[font=\small, text=black!85] (zt) at (3.1, 0, 0)
    {$\mathbf{z}_t$};
\draw[arr] (2.2, 0, 0) -- (2.75, 0, 0);

\draw[arr] (3.45, 0, 0) -- (4.2, 0, 0);

% FC layer (thin tall cuboid)
\cuboid{4.4}{0.25}{1.4}{0.3}{fcfront}{fctop}{fcside}

\draw[arr] (5.0, 0, 0) -- (5.6, 0, 0);

% Q-values output (thin cuboid, lighter)
\cuboid{5.8}{0.25}{0.9}{0.3}{outfront}{outtop}{outside}

% Loss label
\draw[arr] (6.3, 0, 0) -- (7.0, 0, 0);
\node[font=\small, text=palA!80!black] at (7.5, 0, 0)
    {$\mathcal{L}_{\text{Q}}$};

% Component labels (above, same horizontal line)
\node[lbl, anchor=south] at (1.4, 0.85, 0) {encoder};
\node[lbl, anchor=south] at (4.48, 0.85, 0) {Q-head};
\node[lbl, anchor=south] at (5.87, 0.85, 0) {$Q(s,a)$};

% Gradient arrow (below, wider loop)
\draw[gradarr, rounded corners=8pt]
    (1.4, -1.3, 0) -- (1.4, -1.8, 0) -- (7.3, -1.8, 0)
    -- (7.3, -0.3, 0);
\node[lbl, text=palA!50] at (4.3, -2.05, 0)
    {gradient};

\end{tikzpicture}
\caption{The encoder compresses each observation into a
  representation~$\mathbf{z}_t$. The Q-learning loss is the only
  signal that shapes what the encoder learns.}
\label{fig:single-objective}
\end{figure}


\smallskip
The encoder learns by matching its predictions against target values,
but when data is scarce these targets are based on the agent's own
poor estimates. The gradients that update the encoder are small and
inconsistent, so the encoder gets little guidance about what features
to develop.

\smallskip
With limited data, the encoder may find features that satisfy the
current targets without capturing the full structure of the
environment. It might learn that dark pixels at the top of the screen
correlate with high Q-values, because the ball happens to be there in
the small set of stored transitions, rather than learning to track the
ball's position and velocity. The features overfit at the
representation level, working for the data seen so far but failing to
generalize to new states.

\smallskip
With so little training data, good representations must emerge quickly
or not at all. The question is how to provide an additional training
signal that does not depend on the quality of these targets.

\smallskip
One answer is \emph{self-supervised learning}, which creates training
signals from the structure of the data itself, without requiring
external labels. In practice, this means defining a \emph{pretext
task}, an auxiliary objective that requires the encoder to capture
useful information about the input, using only the observations
themselves as supervision.

\smallskip
This idea first proved its value in computer vision, where a common
pretext task is predicting how an image has been rotated. A network
that solves this must learn to recognize the objects and spatial
layout within the image. Modern methods such as
BYOL~\citep{grill2020bootstrap} push this further, learning features
that rival supervised approaches.

\smallskip
Reinforcement learning offers an even richer source of pretext tasks
because the agent observes sequences of states connected by actions,
a property called \emph{temporal structure}. State~$s_t$,
action~$a_t$, and next state~$s_{t+1}$ form a natural prediction
problem. Given the current state and the action taken, predict the
next state.

\smallskip
This prediction task captures \emph{dynamics}, how the game state
evolves under different actions. An encoder that can predict the next
state must represent object positions, velocities, and the effects of
actions. These are exactly the features that are useful for
decision-making.

\smallskip
But predicting the next state in pixel space forces the encoder to
model irrelevant details like exact textures, background patterns, and
rendering artifacts. Instead, the encoder can predict its own future
\emph{representations}, capturing positions, velocities, and how
actions change the game without modeling every pixel. This is the key
idea behind SPR (\cref{subsec:spr}), which trains the encoder to
predict its own future representations, not future images.

\smallskip
This prediction objective is combined with the Q-learning loss:
%
\begin{equation}
\label{eq:ssl-loss}
\mathcal{L}_{\text{total}}
  = \mathcal{L}_{\text{Q}}
  + \lambda\, \mathcal{L}_{\text{SSL}}.
\end{equation}
%
The weight $\lambda$ controls how much the encoder learns from each
loss. The Q-learning loss tells the encoder which features predict
accurate action values. The self-supervised loss tells it how
observations change over time. Together, the Q-learning loss provides
task-relevant guidance even when noisy early in training, and the
self-supervised loss provides guidance about how the game changes even
when it says nothing directly about rewards.

\smallskip
Compared to the single-objective architecture in
\cref{fig:single-objective}, the encoder is now updated by two losses
instead of one (\cref{fig:dual-objective}). The shared encoder now
serves two purposes, computing Q-values and predicting future
representations, and both losses shape what it learns.

\FloatBarrier
\begin{figure}[htbp]
\centering
% Thesis-wide categorical color palette
% Usage: % Thesis-wide categorical color palette
% Usage: \input{figures/palette} in any figure file
%
% Muted primary colors -- print-friendly, visually distinct.
% Inspired by Cooper Union's red/blue primaries, softened for
% academic figures. Use palA-palC for data series in scatter
% plots, line charts, bar charts, etc.

\definecolor{palA}{HTML}{1F77B4}   % Tableau blue
\definecolor{palB}{HTML}{D62728}   % Tableau red
\definecolor{palC}{HTML}{2CA02C}   % Tableau green
 in any figure file
%
% Muted primary colors -- print-friendly, visually distinct.
% Inspired by Cooper Union's red/blue primaries, softened for
% academic figures. Use palA-palC for data series in scatter
% plots, line charts, bar charts, etc.

\definecolor{palA}{HTML}{1F77B4}   % Tableau blue
\definecolor{palB}{HTML}{D62728}   % Tableau red
\definecolor{palC}{HTML}{2CA02C}   % Tableau green

\definecolor{encfade}{HTML}{B8D2E6}
\definecolor{outfront}{HTML}{DCE8F4}
\definecolor{outtop}{HTML}{EBF1F8}
\definecolor{outside}{HTML}{C2D6E8}
\definecolor{fcfront}{HTML}{B3CEE4}
\definecolor{fctop}{HTML}{CDDEF0}
\definecolor{fcside}{HTML}{8DB3D2}
\definecolor{sslfront}{HTML}{B8D8A0}
\definecolor{ssltop}{HTML}{CBE4B8}
\definecolor{sslside}{HTML}{8FB870}

\newcommand{\cuboid}[7]{%
    \fill[#5, draw=black!35, line join=round]
        (#1, {-#3/2}, 0) -- ({#1+#2}, {-#3/2}, 0) --
        ({#1+#2}, {#3/2}, 0) -- (#1, {#3/2}, 0) -- cycle;
    \fill[#6, draw=black!35, line join=round]
        (#1, {#3/2}, 0) -- ({#1+#2}, {#3/2}, 0) --
        ({#1+#2}, {#3/2}, #4) -- (#1, {#3/2}, #4) -- cycle;
    \fill[#7, draw=black!35, line join=round]
        ({#1+#2}, {-#3/2}, 0) -- ({#1+#2}, {#3/2}, 0) --
        ({#1+#2}, {#3/2}, #4) -- ({#1+#2}, {-#3/2}, #4) -- cycle;
}

\begin{tikzpicture}[
    x={(1cm,0cm)}, y={(0cm,1cm)}, z={(0.35cm,0.25cm)},
    arr/.style={->, thick, >=stealth, black!45},
    gradarr/.style={<-, thick, >=stealth, dashed},
    lbl/.style={font=\scriptsize, text=black!45},
]

% Observation label
\node[font=\small, text=black!85] at (-0.5, 0, 0) {$s_t$};
\draw[arr] (0.0, 0, 0) -- (0.6, 0, 0);

% Encoder (trapezoid)
\fill[encfade, draw=black!60, thick, line join=round]
    (0.8, -0.9, 0) -- (0.8, 0.9, 0) -- (2.5, 0.5, 0)
    -- (2.5, -0.5, 0) -- cycle;
\node[font=\small, text=black!80] at (1.55, 0, 0) {$f_\theta$};

% z_t label
\node[font=\small, text=black!85] (zt) at (3.4, 0, 0)
    {$\mathbf{z}_t$};
\draw[arr] (2.5, 0, 0) -- (3.05, 0, 0);

% === Upper path: Q-head ===
\draw[arr] (3.75, 0.15, 0) -- (4.6, 1.8, 0);

% FC layer (upper)
\cuboid{4.8}{0.15}{1.0}{0.1}{fcfront}{fctop}{fcside}
\draw[arr] (5.3, 1.8, 0) -- (6.1, 1.8, 0);

% Q-values output (upper)
\cuboid{6.3}{0.15}{0.6}{0.1}{outfront}{outtop}{outside}
\node[lbl, anchor=south] at (6.37, 2.15, 0) {$Q(s,a)$};

% Q loss
\draw[arr] (6.8, 1.8, 0) -- (7.5, 1.8, 0);
\node[font=\small, text=palA!80!black] at (8.0, 1.8, 0)
    {$\mathcal{L}_{\text{Q}}$};

% Q-head label (between fork and cuboids)
\node[lbl] at (5.05, 1.05, 0) {Q-head};

% === Lower path: SSL head ===
\draw[arr] (3.75, -0.15, 0) -- (4.6, -1.8, 0);

% SSL projection (lower, green cuboid)
\cuboid{4.8}{0.15}{1.0}{0.1}{sslfront}{ssltop}{sslside}
\draw[arr] (5.3, -1.8, 0) -- (6.1, -1.8, 0);

% SSL output (lower, green cuboid)
\cuboid{6.3}{0.15}{0.6}{0.1}{sslfront}{ssltop}{sslside}
\node[lbl, anchor=north] at (6.37, -2.15, 0)
    {$\hat{\mathbf{z}}_{t+1}$};

% SSL loss
\draw[arr] (6.8, -1.8, 0) -- (7.5, -1.8, 0);
\node[font=\small, text=palC!70!black] at (8.2, -1.8, 0)
    {$\lambda\,\mathcal{L}_{\text{SSL}}$};

% SSL-head label
\node[lbl] at (5.05, -1.05, 0) {SSL head};

% Encoder + bottleneck labels
\node[lbl, anchor=north] at (1.55, -1.15, 0) {encoder};
\node[lbl, anchor=north] at (3.4, -0.4, 0) {512};

% === Gradient arrows (colored to show which loss) ===
% Q gradient (blue, above)
\draw[gradarr, palA!55, rounded corners=5pt]
    (1.55, 1.5, 0) -- (1.55, 2.8, 0) -- (7.8, 2.8, 0)
    -- (7.8, 2.05, 0);
\node[lbl, text=palA!50] at (4.7, 3.0, 0) {action values};

% SSL gradient (green, below)
\draw[gradarr, palC!45!black, rounded corners=5pt]
    (1.55, -1.5, 0) -- (1.55, -2.8, 0) -- (8.0, -2.8, 0)
    -- (8.0, -2.05, 0);
\node[lbl, text=palC!45!black] at (4.7, -3.0, 0)
    {dynamics / structure};

\end{tikzpicture}
\caption{Dual-objective training. Self-supervised learning adds a
  second path: the encoder receives gradients from both
  $\mathcal{L}_{\text{Q}}$ (action values) and
  $\lambda\mathcal{L}_{\text{SSL}}$ (dynamics), shaping its
  representations with complementary signals.}
\label{fig:dual-objective}
\end{figure}


\smallskip
\emph{Contrastive methods} such as CURL~\citep{laskin2020curl}
define this self-supervised loss by pulling similar inputs together
and pushing different inputs apart. \emph{Non-contrastive} or
\emph{predictive methods} such as BYOL~\citep{grill2020bootstrap}
and SPR take a different approach, predicting future outputs of a
slowly updated copy of the encoder without requiring negative
examples.

\smallskip
SPR uses the predictive approach, which is a better fit for RL
because it does not require large batches of negative examples. RL
typically trains with batches of 32, far smaller than the thousands
used in vision. Predictive methods are also simpler to implement and
less sensitive to augmentation choices.

\smallskip
But if the encoder maps every input to the same constant vector, the
prediction task is solved perfectly and the resulting representation
is useless. The next section explains this collapse problem and how
modern methods overcome it.

\subsection{Preventing Collapse}
\label{subsec:preventing-collapse}

This failure is called \emph{representational collapse}, and it is
not a rare edge case. It is the global minimum of any naive prediction
loss, and gradient descent will find it unless something prevents it.
The loss only measures how well the prediction matches the target, so
nothing penalizes an encoder that produces constant output. Useful
representations preserve meaningful distinctions between game states,
while collapsed representations map every state to the same point
(\cref{fig:collapse}).

\begin{figure}[htbp]
\centering
% Thesis-wide categorical color palette
% Usage: % Thesis-wide categorical color palette
% Usage: \input{figures/palette} in any figure file
%
% Muted primary colors -- print-friendly, visually distinct.
% Inspired by Cooper Union's red/blue primaries, softened for
% academic figures. Use palA-palC for data series in scatter
% plots, line charts, bar charts, etc.

\definecolor{palA}{HTML}{1F77B4}   % Tableau blue
\definecolor{palB}{HTML}{D62728}   % Tableau red
\definecolor{palC}{HTML}{2CA02C}   % Tableau green
 in any figure file
%
% Muted primary colors -- print-friendly, visually distinct.
% Inspired by Cooper Union's red/blue primaries, softened for
% academic figures. Use palA-palC for data series in scatter
% plots, line charts, bar charts, etc.

\definecolor{palA}{HTML}{1F77B4}   % Tableau blue
\definecolor{palB}{HTML}{D62728}   % Tableau red
\definecolor{palC}{HTML}{2CA02C}   % Tableau green


\begin{subfigure}[t]{0.45\textwidth}
\centering
\begin{tikzpicture}[
    lbl/.style={font=\scriptsize, text=black!55},
]

% Axes
\draw[->, black!30] (-0.2, 0) -- (3.8, 0);
\draw[->, black!30] (0, -0.2) -- (0, 3.5);

% Cluster 1: "ball left" states (palA blue) -- upper left
\fill[palA!70] (0.5, 2.9) circle (3pt);
\fill[palA!70] (0.8, 2.6) circle (3pt);
\fill[palA!70] (0.4, 2.4) circle (3pt);
\fill[palA!70] (1.0, 2.8) circle (3pt);
\fill[palA!70] (0.7, 2.3) circle (3pt);

% Cluster 2: "ball center" states (palC green) -- lower left
\fill[palC!70] (1.0, 0.7) circle (3pt);
\fill[palC!70] (1.3, 1.0) circle (3pt);
\fill[palC!70] (0.8, 1.1) circle (3pt);
\fill[palC!70] (1.2, 0.5) circle (3pt);
\fill[palC!70] (1.5, 0.8) circle (3pt);

% Cluster 3: "ball right" states (palB red) -- right
\fill[palB!70] (3.1, 1.8) circle (3pt);
\fill[palB!70] (3.4, 1.5) circle (3pt);
\fill[palB!70] (2.9, 1.4) circle (3pt);
\fill[palB!70] (3.5, 1.7) circle (3pt);
\fill[palB!70] (3.2, 1.2) circle (3pt);

% Cluster labels -- positioned away from dots
\node[font=\tiny, text=palA!80!black, anchor=east]
    at (0.25, 2.6) {ball left};
\node[font=\tiny, text=palC!80!black, anchor=north]
    at (1.1, 0.25) {ball center};
\node[font=\tiny, text=palB!80!black, anchor=west]
    at (3.7, 1.5) {ball right};

% Prediction arrow example (meaningful)
\draw[->, thick, black!30, shorten >=2pt, shorten <=2pt]
    (1.3, 1.0) -- (2.9, 1.4);
\node[font=\tiny, text=black!40, anchor=south]
    at (2.1, 1.35) {predict};

\end{tikzpicture}
\subcaption{Useful representations}
\label{fig:collapse-useful}
\end{subfigure}
\hfill
\begin{subfigure}[t]{0.45\textwidth}
\centering
\begin{tikzpicture}[
    lbl/.style={font=\scriptsize, text=black!55},
]

% Axes
\draw[->, black!30] (-0.2, 0) -- (3.8, 0);
\draw[->, black!30] (0, -0.2) -- (0, 3.5);

% All points collapsed to center -- spread enough to see colors
\fill[palA!70] (1.75, 1.80) circle (3pt);
\fill[palA!70] (1.95, 1.90) circle (3pt);
\fill[palA!70] (1.65, 1.55) circle (3pt);
\fill[palA!70] (2.00, 1.75) circle (3pt);
\fill[palA!70] (1.80, 1.60) circle (3pt);

\fill[palC!70] (2.05, 1.65) circle (3pt);
\fill[palC!70] (1.85, 1.85) circle (3pt);
\fill[palC!70] (1.60, 1.70) circle (3pt);
\fill[palC!70] (2.10, 1.55) circle (3pt);
\fill[palC!70] (1.70, 1.90) circle (3pt);

\fill[palB!70] (2.05, 1.80) circle (3pt);
\fill[palB!70] (1.75, 1.70) circle (3pt);
\fill[palB!70] (1.90, 1.55) circle (3pt);
\fill[palB!70] (2.00, 1.90) circle (3pt);
\fill[palB!70] (1.65, 1.65) circle (3pt);

% Highlight the collapse cluster
\draw[black!30, thick, dashed]
    (1.85, 1.72) circle (0.40);

% Constant vector label
\node[font=\scriptsize, text=black!50, anchor=north]
    at (1.85, 1.2) {$\mathbf{c}$};

% Annotation: loss = 0, information = 0
\node[font=\scriptsize, text=palB!70!black, align=left,
    anchor=south west] at (2.5, 2.5)
    {loss $= 0$\\[1pt] information $= 0$};
\draw[thin, black!20] (2.25, 1.95) -- (2.5, 2.5);

% Trivial prediction (self-loop, larger)
\draw[->, thick, black!30]
    (2.35, 1.72) arc[start angle=0, end angle=-310,
    x radius=0.5, y radius=0.35];
\node[font=\scriptsize, text=black!40, anchor=west]
    at (2.9, 1.72) {predict $\mathbf{c}$};

\end{tikzpicture}
\subcaption{Collapsed representations}
\label{fig:collapse-collapsed}
\end{subfigure}

\caption{Representational collapse. (a)~The encoder maps different
  game states to distinct regions, preserving meaningful structure.
  (b)~Under collapse, all states map to the same constant
  vector~$\mathbf{c}$: prediction is trivially correct but the
  representation carries no information.}
\label{fig:collapse}
\end{figure}


\smallskip
Contrastive methods such as CURL~\citep{laskin2020curl} prevent
collapse by requiring the encoder to distinguish between different
inputs, but they need large batches of negative examples to work
well. RL typically trains with batches of 32, far too small for
contrastive approaches. BYOL~\citep{grill2020bootstrap} showed that
collapse can be prevented without negative examples using two
mechanisms together. The first, an \emph{exponential moving average
(EMA) target encoder}, is a separate copy whose parameters are a
slowly updated average of the training encoder's parameters:
%
\begin{equation}
\label{eq:ema-update}
\boldsymbol{\theta}_{\text{target}}
  \leftarrow \tau\, \boldsymbol{\theta}_{\text{target}}
  + (1 - \tau)\, \boldsymbol{\theta}_{\text{online}}.
\end{equation}
%
With $\tau$ typically set to 0.99, the target retains 99\% of its
previous parameters at each step and absorbs 1\% from the training
encoder. Over hundreds of steps, the target gradually follows the
training encoder, creating a temporal lag. The target reflects what
the training encoder looked like in the recent past, not what it
looks like now.

\smallskip
This temporal lag makes collapse self-correcting because the target
does not collapse simultaneously. If the training encoder improves,
the target slowly catches up, providing progressively better
prediction targets. If the training encoder tries to collapse, the
target still reflects diverse representations from the recent past,
so the prediction loss increases rather than decreases, pulling the
training encoder back.

\smallskip
The second mechanism, an \emph{asymmetric predictor head}, is a small
network on the training side that maps the encoder's representation
into a form comparable with the target. Because the target side has no
predictor, this mapping must be meaningful. If the encoder collapses,
the predictor receives the same constant input for every example and
cannot learn anything useful, pushing the encoder back toward diverse
representations.

\smallskip
The EMA target encoder and the asymmetric predictor must work
together, and both follow the same principle as the DQN target
network (\cref{subsec:target-networks}), which provides stable
Q-value targets through a frozen copy of the network. The EMA
target encoder serves the same role for representations, but where
DQN copies parameters every $C$ steps, the EMA blends parameters at
every step. In SPR, both mechanisms coexist, with the target network
providing stable Q-value targets and the EMA target encoder providing
stable representation targets, allowing the agent to predict its own
future representations without collapse.

\subsection{Self-Predictive Representations}
\label{subsec:spr}

SPR~\citep{schwarzer2021dataefficient} trains the encoder to
predict its own future representations -- not future pixels, not
future rewards, but future latent states. The idea follows directly
from the previous two subsections: RL has a natural pretext task
(temporal prediction), predicting in latent space avoids the cost of
pixel reconstruction, and the EMA target encoder plus asymmetric
predictor prevent collapse. SPR adds one further insight:
\emph{multi-step} prediction -- predicting $K$ steps ahead rather
than just one -- produces richer representations because the
encoder must capture dynamics that unfold over multiple time steps.
By default, SPR uses $K = 5$.

\smallskip
SPR adds several components on top of the base agent's
architecture. The \emph{online encoder} is the standard CNN
encoder from DQN (the Nature CNN described in
\cref{subsec:dqn-architecture}). It takes the observation $s_t$ and
produces a spatial feature map
$\mathbf{y}_t = f_{\text{online}}(s_t)$. This is the same
encoder that feeds into the Q-head; the SPR components share it
rather than using a separate network.

\smallskip
The \emph{EMA target encoder} is a slow-moving copy of the online
encoder, updated via the EMA rule from \cref{eq:ema-update}. It
takes future observations $s_{t+k}$ and produces target
representations
$\mathbf{y}^{\text{tgt}}_{t+k} = f_{\text{target}}(s_{t+k})$.
Gradients do not flow through this encoder -- it provides stable
targets for the online encoder to predict, using the same
anti-collapse mechanism described in \cref{subsec:preventing-collapse}.

\smallskip
The \emph{transition model} $h$ is an action-conditioned network
that predicts the next latent state given the current latent and
the action taken:
$\hat{\mathbf{y}}_{t+1} = h(\mathbf{y}_t,\, a_t)$. This is the
component that captures dynamics: it learns how the latent
representation should change when the agent takes a specific
action. The transition model is applied iteratively for multi-step
prediction -- each step's output becomes the next step's input.

\smallskip
Two additional heads complete the architecture. A \emph{projection
head} $g$ maps spatial feature maps to a lower-dimensional vector
suitable for comparison. It is applied to both predicted and target
representations. A \emph{prediction head} $q$ -- the asymmetric
predictor from BYOL -- is applied only to the online (predicted)
side, providing the second anti-collapse mechanism from
\cref{subsec:preventing-collapse}. The full comparison for a single
prediction step is between
$q(g(\hat{\mathbf{y}}_{t+k}))$ on the online side and
$g(\mathbf{y}^{\text{tgt}}_{t+k})$ on the target side.
\Cref{fig:spr-architecture} shows the forward prediction loop, and
\cref{fig:spr-loss} shows how the online and target paths are
compared.

\begin{figure}[htbp]
\centering
% Encoder colors (reuse from Nature CNN)
\definecolor{encfade}{HTML}{B8D2E6}
% SPR-specific component color
\definecolor{sprcomp}{HTML}{C2D4A8}
\definecolor{sprfade}{HTML}{D8E4CC}
% Target path (faded)
\definecolor{targetfade}{HTML}{E0E0E0}

\begin{tikzpicture}[
    box/.style={draw=black!40, thick, minimum height=0.75cm,
                font=\small, inner sep=5pt, rounded corners=2pt},
    arr/.style={->, thick, >=stealth, black!40},
    lbl/.style={font=\scriptsize, text=black!55},
    mathlbl/.style={font=\small, text=black!55},
]

% === ONLINE PATH (top) ===
\node[lbl, font=\small\bfseries, text=black!50] at (-0.5, 2.2)
    {Online};

% s_t
\node[mathlbl] at (-0.5, 0.8) {$s_t$};
\draw[arr] (0.0, 0.8) -- (0.8, 0.8);

% Encoder
\fill[encfade, draw=black!30, thick, line join=round]
    (0.8, -0.05) -- (0.8, 1.65) -- (2.3, 1.25) -- (2.3, 0.35) -- cycle;
\node[font=\scriptsize, text=black!45] at (1.45, 0.8) {$f_\theta$};

% z_t
\draw[arr] (2.3, 0.8) -- (3.1, 0.8);
\node[mathlbl] at (2.7, 1.1) {$z_t$};

% Transition model
\node[box, fill=sprcomp] (trans) at (4.0, 0.8) {$h$};
\node[lbl] at (4.0, 0.1) {transition};

% Action input
\node[mathlbl] at (4.0, 2.0) {$a_t$};
\draw[arr] (4.0, 1.7) -- (trans.north);

% z_hat_{t+1}
\draw[arr] (trans.east) -- (5.5, 0.8);
\node[mathlbl] at (5.2, 1.1) {$\hat{z}_{t+1}$};

% Repeated application note
\node[lbl, align=center] at (4.0, -0.5)
    {applied $K$ times};

% Projection
\node[box, fill=sprcomp] (proj) at (6.3, 0.8) {$g$};
\node[lbl] at (6.3, 0.1) {project};

% y_hat
\draw[arr] (proj.east) -- (7.5, 0.8);

% Prediction head
\node[box, fill=sprcomp] (pred) at (8.3, 0.8) {$q$};
\node[lbl] at (8.3, 0.1) {predict};

% pred output
\draw[arr] (pred.east) -- (9.5, 0.8);
\node[mathlbl] at (9.9, 0.8) {$\tilde{y}$};

% === TARGET PATH (bottom) ===
\node[lbl, font=\small\bfseries, text=black!50] at (-0.5, -2.8)
    {Target};

% s_{t+k}
\node[mathlbl] at (-0.5, -3.7) {$s_{t\!+\!k}$};
\draw[arr, black!25] (0.1, -3.7) -- (0.8, -3.7);

% EMA Encoder
\fill[targetfade, draw=black!20, thick, line join=round]
    (0.8, -4.55) -- (0.8, -2.85) -- (2.3, -3.25) -- (2.3, -4.15) -- cycle;
\node[font=\scriptsize, text=black!35] at (1.45, -3.7) {$\bar{f}_\theta$};
\node[lbl, text=black!30] at (1.55, -4.85) {EMA};

% z_tilde
\draw[arr, black!25] (2.3, -3.7) -- (3.1, -3.7);

% EMA Projection
\node[box, fill=targetfade, draw=black!20, text=black!45]
    (tproj) at (3.9, -3.7) {$\bar{g}$};
\node[lbl, text=black!30] at (3.9, -4.4) {EMA};

% target output
\draw[arr, black!25] (tproj.east) -- (5.5, -3.7);
\node[mathlbl, text=black!45] at (5.9, -3.7) {$\bar{y}$};

% === EMA UPDATE ARROW ===
\draw[->, thick, dashed, black!30]
    (1.55, -0.05) -- (1.55, -2.85)
    node[midway, right=3pt, lbl, text=black!30] {EMA};

% === COSINE SIMILARITY ===
\draw[<->, thick, black!50]
    (9.9, 0.4) -- (9.9, -1.5)
    node[midway, right=3pt, font=\small, text=black!50]
    {$\mathcal{L}_{\mathrm{SPR}}$};

% Connect target output to loss
\draw[->, thick, dashed, black!30]
    (5.9, -3.3) -- (9.5, -1.5);

% Gradient annotation
\draw[decorate, decoration={brace, amplitude=4pt}, black!30]
    (0.7, 2.0) -- (9.0, 2.0)
    node[midway, above=5pt, lbl, text=black!35]
    {gradients flow (online path only)};

\end{tikzpicture}
\caption{SPR training architecture. The online encoder $f_\theta$
  produces spatial features $z_t$, which the transition model $h$
  rolls forward $K$ steps using actions. Projected predictions are
  compared via cosine similarity to targets from the EMA encoder
  $\bar{f}_\theta$. The prediction head $q$ exists only on the
  online side, creating the asymmetry that prevents representation
  collapse. Gradients flow through the online path only; the target
  network is updated via exponential moving average.}
\label{fig:spr-architecture}
\end{figure}

\begin{figure}[htbp]
\centering
% Encoder colors (reuse from Nature CNN)
\definecolor{encfade}{HTML}{B8D2E6}
% SPR-specific component color
\definecolor{sprcomp}{HTML}{C2D4A8}
% Target path (faded)
\definecolor{targetfade}{HTML}{E0E0E0}

\begin{tikzpicture}[
    box/.style={draw=black!70, thick, minimum height=0.7cm,
                font=\small, inner sep=5pt, rounded corners=2pt},
    arr/.style={->, thick, >=stealth, black!55},
    lbl/.style={font=\scriptsize, text=black!55},
]

% === Online path (top) ===
\node[font=\scriptsize\bfseries, text=black!70, anchor=west] at (-0.7, 0.7) {Online};

% z_hat_{t+k} input
\node[font=\small, text=black!90] at (-0.3, 0) {$\hat{z}_{t\!+\!k}$};
\draw[arr] (0.3, 0) -- (1.0, 0);

% Projection
\node[box, fill=sprcomp] (proj) at (1.8, 0) {$g$};
\node[lbl] at (1.8, -0.7) {project};

\draw[arr] (proj.east) -- (3.2, 0);

% Prediction head
\node[box, fill=sprcomp] (pred) at (4.0, 0) {$q$};
\node[lbl] at (4.0, -0.7) {predict};

% Online output -- align at x=5.65 for clean vertical comparison
\draw[arr] (pred.east) -- (5.3, 0);
\node[font=\small, text=black!90] (ytilde) at (5.65, 0) {$\tilde{y}$};

% === Target path (bottom) ===
\node[font=\scriptsize\bfseries, text=black!70, anchor=west] at (-0.7, -2.5) {Target};

% s_{t+k}
\node[font=\small, text=black!70] at (-0.5, -3.2) {$s_{t\!+\!k}$};
\draw[arr, black!50] (0.1, -3.2) -- (0.6, -3.2);

% EMA Encoder (trapezoid)
\fill[targetfade, draw=black!55, thick, line join=round]
    (0.8, -4.0) -- (0.8, -2.4) -- (2.1, -2.75) -- (2.1, -3.65) -- cycle;
\node[font=\small, text=black!70] at (1.35, -3.2) {$\bar{f}_\theta$};

\draw[arr, black!50] (2.1, -3.2) -- (2.8, -3.2);

% EMA Projection
\node[box, fill=targetfade, draw=black!55, text=black!70]
    (tproj) at (3.6, -3.2) {$\bar{g}$};

% Target output -- align at x=5.65 to match online output
\draw[arr, black!50] (tproj.east) -- (5.3, -3.2);
\node[font=\small, text=black!70] (ybar) at (5.65, -3.2) {$\bar{y}$};

% === Cosine similarity ===
% Single vertical dashed line connecting y_tilde and y_bar
\draw[thick, dashed, black!55]
    (5.65, -0.3) -- (5.65, -2.9);
\node[font=\small, text=black!85, anchor=west] at (6.3, -1.6)
    {$\mathcal{L}_{\text{SPR}} = -\cos(\tilde{y},\, \bar{y})$};

% No prediction head annotation
\draw[decorate, decoration={brace, amplitude=3pt, mirror}, black!45]
    (2.5, -4.1) -- (4.2, -4.1)
    node[midway, below=4pt, font=\scriptsize, text=black!55]
    {no prediction head (asymmetry)};

% EMA update arrow (between the two paths)
\node[font=\scriptsize, text=black!55, align=center]
    at (-0.5, -1.5) {EMA\\update};
\draw[->, thick, dashed, black!45]
    (-0.5, -0.5) -- (-0.5, -1.1);
\draw[->, thick, dashed, black!45]
    (-0.5, -1.9) -- (-0.5, -2.3);

\end{tikzpicture}
\caption{Self-predictive loss. The predicted representation
  $\tilde{y}$ from the online path is compared to the target
  $\bar{y}$ from an EMA copy of the encoder via cosine similarity.
  The prediction head $q$ appears only on the online side, creating
  the asymmetry that prevents both paths from collapsing to a
  trivial constant.}
\label{fig:spr-loss}
\end{figure}


\smallskip
SPR's training operates on sequences sampled from the replay
buffer. Given a transition sequence
$(s_t, a_t, s_{t+1}, a_{t+1}, \ldots, s_{t+K})$, the procedure
is as follows. First, the current state is encoded with the
online encoder:
$\mathbf{y}_t = f_{\text{online}}(s_t)$. Then, for each of the
next $K$ steps, the transition model predicts the next latent from
the previous \emph{prediction} (not from a fresh encoding):
$\hat{\mathbf{y}}_{t+k} = h(\hat{\mathbf{y}}_{t+k-1},\, a_{t+k-1})$.
Meanwhile, the actual future state $s_{t+k}$ is encoded with the
target encoder to produce the prediction target
$\mathbf{y}^{\text{tgt}}_{t+k}$. The SPR loss is the negative mean
cosine similarity across all $K$ steps:
\begin{equation}
\label{eq:spr-loss}
\mathcal{L}_{\text{SPR}}
  = -\frac{1}{K} \sum_{k=1}^{K}
    \frac{q(g(\hat{\mathbf{y}}_{t+k})) \cdot
          g(\mathbf{y}^{\text{tgt}}_{t+k})}
         {\|q(g(\hat{\mathbf{y}}_{t+k}))\|\;
          \|g(\mathbf{y}^{\text{tgt}}_{t+k})\|}.
\end{equation}
Gradients flow through the online encoder, transition model,
projection head, and prediction head, but not through the target
encoder.

\smallskip
The chained predictions are a deliberate design choice, not a
shortcut. Because each step uses the \emph{previous prediction}
rather than a fresh encoding, errors compound: a poor prediction at
step~1 produces a flawed input for step~2, which in turn degrades
step~3, and so on. This compounding pressure forces the encoder to
produce representations that the transition model can propagate
accurately over multiple steps. The encoder must represent
information that is \emph{predictable over time} -- stable features
such as object positions and velocities, rather than transient
noise like exact pixel patterns. The $K = 5$ horizon is enough to
capture short-term dynamics (ball trajectories, paddle movements)
without requiring long-horizon planning.

\smallskip
The encoder thus receives gradients from two losses: the Q-learning
loss, which signals what actions are valuable, and the SPR loss,
which signals how states evolve under actions. The SPR loss
encourages the encoder to represent object positions, velocities,
and action effects -- information that is useful for Q-learning but
that the Q-loss alone may not explicitly require, especially early
in training when Q-targets are noisy. The result is richer, more
informative representations that support faster and more stable
Q-learning. The encoder captures the environment's dynamics, not
just whatever happens to correlate with current Q-targets.

\smallskip
SPR shapes representations but does not change the agent's
exploration strategy, replay sampling, or learning algorithm. The
agent still uses $\varepsilon$-greedy exploration and uniform
replay. SPR does not add domain knowledge: the transition model is
a generic convolutional network, and the representations it
encourages are discovered from data. Its benefits may therefore
depend on the base agent. A capable base agent such as Rainbow
provides richer replay data (through prioritized sampling and
$n$-step returns) and more effective exploration (through noisy
networks), which could make SPR's predictions more informative. A
simple DQN may provide weaker data for SPR to learn from.

\smallskip
This is precisely the question that motivates the experiments in
this work. SPR was originally evaluated on top of
Rainbow~\citep{hessel2018rainbow}, which includes six extensions
beyond DQN. The reported improvements on Atari-100K were
substantial, but the base agent was already sophisticated. The
gains could come from SPR's representation learning alone, from
the interaction between SPR and Rainbow's components, or from a
combination of both. By testing SPR on a simple DQN -- without
Rainbow's extensions -- the contribution of representation learning
can be isolated from the contribution of the base agent. Combined
with data augmentation as a second experimental factor, this
produces a controlled study of when self-supervised representations
help and what determines their usefulness.
